\begin{thebibliography}{99}

\bibitem[Alb34]{Alb34}
A. A. Albert,
\emph{On the construction of Riemann matrices. I},
Ann. of Math. (2) \textbf{35} (1934), no.~1, 1--28.
\href{https://doi.org/10.2307/1968116}{\texttt{doi:10.2307/1968116}}.

\bibitem[Alb35a]{Alb35a}
A. A. Albert,
\emph{On the construction of Riemann matrices. II},
Ann. of Math. (2) \textbf{36} (1935), no.~2, 376--394.
\href{https://doi.org/10.2307/1968578}{\texttt{doi:10.2307/1968578}}.

\bibitem[Alb35b]{Alb35b}
A. A. Albert,
\emph{Involutorial simple algebras and real Riemann matrices},
Ann. of Math. (2) \textbf{36} (1935), no.~4, 886--964.
\href{https://doi.org/10.2307/1968595}{\texttt{doi:10.2307/1968595}}.

\bibitem[And96a]{And96}
Y. Andr\'e,
\emph{On the Shafarevich and Tate conjectures for hyperk\"ahler varieties},
Math. Ann. \textbf{305} (1996), no.~1, 205--248.
\href{https://doi.org/10.1007/BF01444219}{\texttt{doi:10.1007/BF01444219}}.

\bibitem[BB26]{BB26}
D. Bhattacharjee and U. Bhattacharya,
\emph{Relative secant cycles and Hodge classes},
Preprints.org, version 5, 2026.
\href{https://doi.org/10.20944/preprints202602.0462.v5}{\texttt{doi:10.20944/preprints202602.0462.v5}}.

\bibitem[BB66]{BB66}
W. L. Baily, Jr. and A. Borel,
\emph{Compactification of arithmetic quotients of bounded symmetric domains},
Ann. of Math. (2) \textbf{84} (1966), no.~3, 442--528.
\href{https://doi.org/10.2307/1970457}{\texttt{doi:10.2307/1970457}}.

\bibitem[BF03]{BF03}
R.-O. Buchweitz and H. Flenner,
\emph{A semiregularity map for modules and applications to deformations},
Compositio Math. \textbf{137} (2003), no.~2, 135--210.
\href{https://doi.org/10.1023/A:1023999012081}{\texttt{doi:10.1023/A:1023999012081}}.

\bibitem[BF08]{BF08}
R.-O. Buchweitz and H. Flenner,
\emph{The global decomposition theorem for Hochschild (co-)homology of
singular spaces via the Atiyah--Chern character},
Adv. Math. \textbf{217} (2008), no.~1, 243--281.
\href{https://doi.org/10.1016/j.aim.2007.06.013}{\texttt{doi:10.1016/j.aim.2007.06.013}}.

\bibitem[BL04]{BL04}
C. Birkenhake and H. Lange,
\emph{Complex Abelian Varieties}, 2nd ed.,
Grundlehren der mathematischen Wissenschaften \textbf{302},
Springer, Berlin, 2004.
\href{https://doi.org/10.1007/978-3-662-06307-1}{\texttt{doi:10.1007/978-3-662-06307-1}}.

\bibitem[Blo72]{Blo72}
S. Bloch,
\emph{Semi-regularity and de Rham cohomology},
Invent. Math. \textbf{17} (1972), 51--66.
\href{https://doi.org/10.1007/BF01390023}{\texttt{doi:10.1007/BF01390023}}.

\bibitem[BMB26]{BMB26}
D. Bhattacharjee, P. Mandal and U. Bhattacharya,
\emph{Verification code and figure sources for the paper ``Algebraic Loci of
Weil Classes on Abelian Varieties''}, software archive with a Lean~4
certificate, release 5.0.1, Zenodo, 2026.
\href{https://doi.org/\zenodoCodeID}{\texttt{doi:\zenodoCodeID}}.

\bibitem[Cal05]{Cal05}
A. C\u{a}ld\u{a}raru,
\emph{The Mukai pairing, II: the Hochschild--Kostant--Rosenberg isomorphism},
Adv. Math. \textbf{194} (2005), no.~1, 34--66.
\href{https://doi.org/10.1016/j.aim.2004.05.012}{\texttt{doi:10.1016/j.aim.2004.05.012}}.

\bibitem[CDK95]{CDK95}
E. Cattani, P. Deligne and A. Kaplan,
\emph{On the locus of Hodge classes},
J. Amer. Math. Soc. \textbf{8} (1995), no.~2, 483--506.
\href{https://doi.org/10.1090/S0894-0347-1995-1273413-2}{\texttt{doi:10.1090/S0894-0347-1995-1273413-2}}.

\bibitem[CKS86]{CKS86}
E. Cattani, A. Kaplan and W. Schmid,
\emph{Degeneration of Hodge structures},
Ann. of Math. (2) \textbf{123} (1986), no.~3, 457--535.
\href{https://doi.org/10.2307/1971333}{\texttt{doi:10.2307/1971333}}.

\bibitem[Del72]{Del72}
P. Deligne,
\emph{La conjecture de Weil pour les surfaces $K3$},
Invent. Math. \textbf{15} (1972), 206--226.
\href{https://doi.org/10.1007/BF01404126}{\texttt{doi:10.1007/BF01404126}}.

\bibitem[Del82]{Del82}
P. Deligne,
\emph{Hodge cycles on abelian varieties} (notes by J. S. Milne),
in: Hodge Cycles, Motives, and Shimura Varieties,
Lecture Notes in Math. \textbf{900}, Springer, Berlin, 1982, pp.~9--100.
\href{https://doi.org/10.1007/978-3-540-38955-2_3}%
{\texttt{doi:10.1007/978-3-540-38955-2\_3}}.

\bibitem[FF26]{FF26}
S. Floccari and L. Fu,
\emph{The Hodge conjecture for Weil fourfolds with discriminant $1$ via
singular OG6-varieties},
J. Math. Pures Appl. (9) \textbf{210} (2026), paper no.~103876.
\href{https://doi.org/10.1016/j.matpur.2026.103876}{\texttt{doi:10.1016/j.matpur.2026.103876}}.

\bibitem[Ful98]{Ful98}
W. Fulton,
\emph{Intersection Theory}, 2nd ed.,
Ergebnisse der Mathematik und ihrer Grenzgebiete (3) \textbf{2},
Springer, New York, 1998.
\href{https://doi.org/10.1007/978-1-4612-1700-8}{\texttt{doi:10.1007/978-1-4612-1700-8}}.

\bibitem[GP98]{GP98}
M. Gross and S. Popescu,
\emph{Equations of $(1,d)$-polarized abelian surfaces},
Math. Ann. \textbf{310} (1998), no.~2, 333--377.
\href{https://doi.org/10.1007/s002080050151}{\texttt{doi:10.1007/s002080050151}}.

\bibitem[Gro61]{Gro61}
A. Grothendieck,
\emph{Techniques de construction et th\'eor\`emes d'existence en g\'eom\'etrie
alg\'ebrique IV: les sch\'emas de Hilbert},
S\'eminaire Bourbaki, Vol.~6 (1960/61), Exp. No.~221, 249--276;
reprinted by Soc. Math. France, Paris, 1995.

\bibitem[HM73]{HM73}
G. Horrocks and D. Mumford,
\emph{A rank $2$ vector bundle on $\mathbf{P}^{4}$ with $15{,}000$ symmetries},
Topology \textbf{12} (1973), 63--81.
\href{https://doi.org/10.1016/0040-9383(73)90022-0}{\texttt{doi:10.1016/0040-9383(73)90022-0}}.

\bibitem[Kle68]{Kle68}
S. L. Kleiman,
\emph{Algebraic cycles and the Weil conjectures},
in: Dix expos\'es sur la cohomologie des sch\'emas,
North-Holland, Amsterdam, 1968, pp.~359--386.

\bibitem[Kle94]{Kle94}
S. L. Kleiman,
\emph{The standard conjectures},
in: Motives (Seattle, WA, 1991),
Proc. Sympos. Pure Math. \textbf{55}, Part 1,
Amer. Math. Soc., Providence, RI, 1994, pp.~3--20.
\href{https://doi.org/10.1090/pspum/055.1/1265519}{\texttt{doi:10.1090/pspum/055.1/1265519}}.

\bibitem[KOU23]{KOU}
B. Klingler, A. Otwinowska and D. Urbanik,
\emph{On the fields of definition of Hodge loci},
Ann. Sci. \'Ec. Norm. Sup\'er. (4) \textbf{56} (2023), no.~4, 1299--1312.
\href{https://doi.org/10.24033/asens.2555}{\texttt{doi:10.24033/asens.2555}}.

\bibitem[KS67]{KS67}
M. Kuga and I. Satake,
\emph{Abelian varieties attached to polarized $K3$-surfaces},
Math. Ann. \textbf{169} (1967), 239--242.
\href{https://doi.org/10.1007/BF01399540}{\texttt{doi:10.1007/BF01399540}}.

\bibitem[Mar25a]{Mar25a}
E. Markman,
\emph{Cycles on abelian $2n$-folds of Weil type from secant sheaves on abelian
$n$-folds}, preprint, arXiv:2502.03415 (2025).
\href{https://doi.org/10.48550/arXiv.2502.03415}{\texttt{doi:10.48550/arXiv.2502.03415}}.

\bibitem[Mar25b]{Mar25b}
E. Markman,
\emph{Secant sheaves and Weil classes on abelian varieties},
preprint, arXiv:2509.23403 (2025).
\href{https://doi.org/10.48550/arXiv.2509.23403}{\texttt{doi:10.48550/arXiv.2509.23403}}.

\bibitem[Mar25c]{Mar25c}
E. Markman,
\emph{Secant sheaves on abelian $n$-folds with real multiplication and Weil
classes on abelian $2n$-folds with complex multiplication},
preprint, 2025.
\href{https://arxiv.org/abs/2509.23079}{\texttt{arXiv:2509.23079}}.

\bibitem[Muk81]{Muk81}
S. Mukai,
\emph{Duality between $D(X)$ and $D(\hat{X})$ with its application to Picard
sheaves},
Nagoya Math. J. \textbf{81} (1981), 153--175.
\href{https://doi.org/10.1017/S002776300001922X}{\texttt{doi:10.1017/S002776300001922X}}.

\bibitem[MZ95]{MZ95}
B. Moonen and Yu. G. Zarhin,
\emph{Hodge classes and Tate classes on simple abelian fourfolds},
Duke Math. J. \textbf{77} (1995), no.~3, 553--581.
\href{https://doi.org/10.1215/S0012-7094-95-07717-5}{\texttt{doi:10.1215/S0012-7094-95-07717-5}}.

\bibitem[MZ99]{MZ99}
B. Moonen and Yu. G. Zarhin,
\emph{Hodge classes on abelian varieties of low dimension},
Math. Ann. \textbf{315} (1999), no.~4, 711--733.
\href{https://doi.org/10.1007/s002080050333}{\texttt{doi:10.1007/s002080050333}}.

\bibitem[II]{paperII}
D. Bhattacharjee and U. Bhattacharya,
\emph{Secant objects for Weil classes on abelian varieties},
companion to the present paper, preprint (2026).

\bibitem[III]{paperIII}
D. Bhattacharjee and U. Bhattacharya,
\emph{Hodge classes beyond the Weil lines},
companion to the present paper, preprint (2026).

\bibitem[PR94]{PR94}
V. P. Platonov and A. S. Rapinchuk,
\emph{Algebraic Groups and Number Theory},
Pure and Applied Mathematics \textbf{139},
Academic Press, Boston, MA, 1994.

\bibitem[Pri24]{Pri24}
J. P. Pridham,
\emph{Semiregularity as a consequence of Goodwillie's theorem},
Forum Math. Sigma \textbf{12} (2024), e126. arXiv:1208.3111.
\href{https://doi.org/10.1017/fms.2024.132}{\texttt{doi:10.1017/fms.2024.132}}.

\bibitem[PST21]{PST}
J. Pila, A. N. Shankar and J. Tsimerman,
\emph{Canonical heights on Shimura varieties and the Andr\'e--Oort
conjecture}, with an appendix by H. Esnault and M. Groechenig,
preprint, arXiv:2109.08788 (2021).
\href{https://doi.org/10.48550/arXiv.2109.08788}{\texttt{doi:10.48550/arXiv.2109.08788}}.

\bibitem[Sat66]{Sat66}
I. Satake,
\emph{Clifford algebras and families of abelian varieties},
Nagoya Math. J. \textbf{27} (1966), 435--446.
\href{https://doi.org/10.1017/S0027763000026295}{\texttt{doi:10.1017/S0027763000026295}}.

\bibitem[Sch73]{Sch73}
W. Schmid,
\emph{Variation of Hodge structure: the singularities of the period mapping},
Invent. Math. \textbf{22} (1973), 211--319.
\href{https://doi.org/10.1007/BF01389674}{\texttt{doi:10.1007/BF01389674}}.

\bibitem[Sch85]{Sch85}
W. Scharlau,
\emph{Quadratic and Hermitian Forms},
Grundlehren der mathematischen Wissenschaften \textbf{270},
Springer, Berlin, 1985.
\href{https://doi.org/10.1007/978-3-642-69971-9}{\texttt{doi:10.1007/978-3-642-69971-9}}.

\bibitem[Sch98]{Sch98}
C. Schoen,
\emph{Addendum to: Hodge classes on self-products of a variety with an
automorphism},
Compositio Math. \textbf{114} (1998), no.~3, 329--336.
\href{https://doi.org/10.1023/A:1000566205021}{\texttt{doi:10.1023/A:1000566205021}}.

\bibitem[Tot26]{Tot26}
B. Totaro,
\emph{On the Hodge conjecture},
guest post on the weblog of T. Tao, 11 September 2026.
\href{https://terrytao.wordpress.com/2026/09/11/on-the-hodge-conjecture/}%
{\texttt{terrytao.wordpress.com/2026/09/11/on-the-hodge-conjecture}}.

\bibitem[Tsi18]{Tsi18}
J. Tsimerman,
\emph{The Andr\'e--Oort conjecture for $\mathcal{A}_{g}$},
Ann. of Math. (2) \textbf{187} (2018), no.~2, 379--390.
\href{https://doi.org/10.4007/annals.2018.187.2.2}{\texttt{doi:10.4007/annals.2018.187.2.2}}.

\bibitem[vG00]{vG00}
B. van Geemen,
\emph{Kuga--Satake varieties and the Hodge conjecture},
in: The Arithmetic and Geometry of Algebraic Cycles (Banff, 1998),
NATO Sci. Ser. C Math. Phys. Sci. \textbf{548},
Kluwer, Dordrecht, 2000, pp.~51--82.
\href{https://doi.org/10.1007/978-94-011-4098-0_3}{\texttt{doi:10.1007/978-94-011-4098-0\_3}}.

\bibitem[Voi02]{Voi02}
C. Voisin,
\emph{Hodge Theory and Complex Algebraic Geometry I},
Cambridge Studies in Advanced Mathematics \textbf{76},
Cambridge University Press, Cambridge, 2002.
\href{https://doi.org/10.1017/CBO9780511615344}{\texttt{doi:10.1017/CBO9780511615344}}.

\bibitem[Voi03]{Voi03}
C. Voisin,
\emph{Hodge Theory and Complex Algebraic Geometry II},
Cambridge Studies in Advanced Mathematics \textbf{77},
Cambridge University Press, Cambridge, 2003.
\href{https://doi.org/10.1017/CBO9780511615177}{\texttt{doi:10.1017/CBO9780511615177}}.

\bibitem[Voi26]{Voi26}
C. Voisin,
\emph{The status of the Hodge conjecture},
guest post on the weblog of T. Tao, 12 September 2026.
\href{https://terrytao.wordpress.com/2026/09/12/the-status-of-the-hodge-conjecture/}%
{\texttt{terrytao.wordpress.com/2026/09/12/the-status-of-the-hodge-conjecture}}.

\bibitem[Wei77]{Weil77}
A. Weil,
\emph{Abelian varieties and the Hodge ring},
in: \OE uvres Scientifiques, Collected Papers, Vol.~III,
Springer, New York, 1979, pp.~421--429.

\end{thebibliography}
